% !TeX root = ../Book1.tex
\begin{figure}[htbp]
	\centering
	\begin{tikzpicture}[font=\small,scale=1.05]
		\draw[->,GeometryGray] (-1,0)--(4.2,0) node[right] {\(E\)};
		\draw[->,GeometryGray] (0,-0.65)--(0,3.25) node[above] {\(Q\)};
		\draw[GeometryBlue,very thick] (-0.7,-0.455)--(4,2.6) node[above] {\(W_A\)};
		\draw[->,thick] (0,0)--(2.8,0) node[midway,below] {\(x\)};
		\draw[->,orange!75!black,thick] (2.8,0)--(2.8,1.82) node[midway,right] {\(Ax\)};
		\draw[->,GeometryBlue,thick] (0,0)--(2.8,1.82);
		\fill[GeometryBlue] (2.8,1.82) circle[radius=2pt];
		\node[above left] at (2.8,1.82) {\(x+Ax\)};
		\node[below left] at (0,0) {\(0\)};
		\node[align=center] at (6.5,1.65) {投影 \(W_A\to E\) 可逆\\[4pt]每个 \(x\) 唯一补上 \(Ax\)};
	\end{tikzpicture}
	\caption[Grassmann 流形的图坐标]{图坐标以线性映射 \(A:E\to Q\) 记录子空间 \(W_A\)。图画的是一维横向与一维竖向情形；投影可逆是使用这张坐标图的条件。}
	\label{g1:fig:grassmann-graph}
\end{figure}
